\section{Exact affine event discovery and atomic hinges}
\label{sec:r16-events}

The R16 profile \code{ATOMOS-EXACT-AFFINE-EVENTS-R1} specializes the
source world-query continuation idea to rational affine trajectories
and the existing editable aTOMos plane guards. The implementation is
\path{python/atomos_events.py}, using
\code{Engine.step\_certified\_crossing} in
\path{python/atomos_hinge.py}. It is an exact host event adapter, distinct
from the native finite-word VM in Section~\ref{sec:r16-wqk}. It does
not claim GPU analytic root discovery or dynamic GPU spatial refitting.

The reviewed source basis is the WQK package's
\path{spec/TOM_WORLD_QUERY_KERNEL_0_4_REBUILT.md}, continuation schema,
and \path{src/python/tom_world04r/} implementation. R16 retains the
useful event/continuation structure while specifying strict rational
admission, complete semantic replay, simultaneous conflict handling,
and terminal-horizon behavior.

\subsection{One plane definition and one immutable epoch}

An epoch $E$ contains rational start $t_0$, finite rational horizon
$H>t_0$, state $x_0$, rates $v$, a finite set of unconsumed relations,
and fixed nominal frame and time convention. Each scalar field declares
seven integer SI dimension exponents; a rate has that field's dimension
divided by time. Between events,
\begin{equation}
 x(t)=x_0+v(t-t_0),\qquad t\in[t_0,H].
 \label{eq:r16-affine-epoch}
\end{equation}
This equation is the current epoch's conditional local prediction;
an earlier accepted event replaces its remaining continuation. It does
not assert that an arbitrary accelerating physical body is exactly affine.

A relation selects three length fields and binds them to the existing
literal plane expression in a hinge seed. With rational normal $n$
and offset $d$, the guard and its pullback are
\begin{align}
 g(x)&=n\cdot x-d,\nonumber\\
 g(x(t))&=a(t-t_0)+c,
 &a&=n\cdot v,&c&=n\cdot x_0-d.
 \label{eq:r16-affine-guard}
\end{align}
The adapter derives $n,d$ from the actual plane graph; the relation
does not carry an independent second guard equation. The field names,
seed identity, and nominal frame are part of the binding. A nonunit
normal is valid for its zero set and crossing signs, but $g$ is then
a defining function rather than a Euclidean signed distance. Unsupported
nonlinear, spherical, or otherwise unrecognized guard graphs reject
this profile rather than falling back to approximate root sampling.

Exact input means integers or reduced rational values. Serialized
\code{num}/\code{den} limbs must be actual integers with positive
denominator. Floating values, Boolean integer aliases, and lossy
integer coercion are excluded. Rational reduction changes representation,
not the represented value.

\subsection{Complete earliest roots and exact exclusions}

Each remaining relation $r$ declares a closed active interval,
inclusive scalar support intervals, and scalar equalities. Its root
search window is
\begin{equation}
 W_r=(t_0,H]\cap[\alpha_r,\beta_r].
 \label{eq:r16-event-window}
\end{equation}
Thus a root at the current start is excluded, including a zero created
by an atomic reset at that timestamp. There is no implicit zero-time
event cascade.

When $a_r\ne0$, the sole root is
\begin{equation}
 \tau_r=t_0-\frac{c_r}{a_r}.
 \label{eq:r16-affine-root}
\end{equation}
It is admitted only if $\tau_r\in W_r$ and every support/equality gate
holds exactly at $x(\tau_r)$. When $a_r=0,c_r\ne0$, there is no root.
When $a_r=c_r=0$, the guard is identically zero and has no isolated
crossing; the result remains unresolved unless a time or gate exclusion
already proves the relation irrelevant over the whole forward window.

For an affine scalar field $x_j$, its exact range on a closed enclosing
interval $[l,u]$ is
\begin{equation}
 I_j=[\min(x_j(l),x_j(u)),\max(x_j(l),x_j(u))].
 \label{eq:r16-affine-enclosure}
\end{equation}
If $I_j$ is disjoint from a required support interval, or does not contain
a required equality value, that gate excludes the entire window. Including
the open window's lower endpoint in this enclosure is conservative.
Failure of this exclusion test does not assert that all gates can hold
simultaneously; it preserves a possible relation. For a nonconstant
guard, the subsequent exact gate test at its unique root resolves it.

Let $\mathcal C(E)$ index the relations with admitted roots, after
checking every unconsumed relation. If nonempty, select
\begin{equation}
 \tau_* = \min_{r\in\mathcal C(E)}\tau_r,\qquad
 \mathcal E_* =\{r\in\mathcal C(E):\tau_r=\tau_*\}.
 \label{eq:r16-earliest-set}
\end{equation}
Priority and relation identity order the certificate within these exact
ties. They neither move the root nor authorize partial chronological
execution inside a simultaneous set. A certificate records all candidate
roots as well as the complete selected earliest set. An unresolved
eligible relation prevents a false completeness claim.

\begin{proposition}[Earliest-event completeness in the affine profile]
If every eligible relation is either exactly resolved or excluded by
the stated interval proof, Equation~\eqref{eq:r16-earliest-set} contains
every admissible event at the least admissible time in $(t_0,H]$.
\end{proposition}
\begin{proof}
A nonconstant affine guard has exactly the root in
Equation~\eqref{eq:r16-affine-root}. A constant nonzero guard has none.
An excluded relation cannot satisfy a required gate in its window.
Every remaining isolated root is checked for window membership and
all gates. Enumeration therefore neither omits an admissible root nor
adds an inadmissible one. Exact rational minimum and equality then
select precisely all least-time roots. The assertion covers the
declared finite relation set and affine algebra only.
\end{proof}

\subsection{A crossing certificate without an epsilon step}

At an admitted nonzero-slope root,
\begin{equation}
 g(x(\tau_*+s))=a s,\qquad
 \sigma_-=-\operatorname{sgn}(a),\quad
 \sigma_+=\operatorname{sgn}(a).
 \label{eq:r16-crossing-sides}
\end{equation}
These signs follow algebraically for $s<0$ and $s>0$; the runtime does
not choose a numerical epsilon. The hinge adapter re-derives slope,
root, point, and source seed binding before using the certified sides.
Its existing chart orientation, if present in the guard convention,
corrects both signs consistently; this fixed-frame crossing itself
declares no reversing seam.

The post-side inside bit drives the selected lane through the actual
two stored mask stages and the existing word $J,K$ bodies. The update
reads an old hinge snapshot and commits once, with event parity toggled
once. Other lanes retain their declared bodies; the optional seam lane
does not receive a toggle from this fixed-frame event.

If several simultaneous relations name the same hinge, they generate
one pulse only when their complete origin, velocity, root, and plane-seed
bindings agree. Different bindings to that hinge are a conflict. Merely
sharing a timestamp is insufficient to identify two physical crossings.

At $\tau_*=H$, $\sigma_+$ is the sign of the algebraic affine extension.
It does not claim an observed or evolved future interval. The horizon
event yields a \code{terminal\_state} point snapshot; it does not construct
a zero-width successor affine segment. With no admitted events, the
adapter advances the affine state to $H$ and finalizes without a hinge
pulse.

\subsection{Atomic simultaneous state and rate updates}

Let $x^-=x(\tau_*)$ and $v^-$ denote the common pre-event snapshot.
For a field with simultaneous updates, the supported rules are
\begin{equation}
 z^+=\begin{cases}
 c,&\text{all updates are set to the same value }c,\\
 z^-+\sum_j\delta_j,&\text{all updates are add},\\
 z^-\oplus\delta_1\oplus\cdots\oplus\delta_m,
      &\text{all updates are admissible xor}.
 \end{cases}
 \label{eq:r16-simultaneous-updates}
\end{equation}
XOR is limited to nonnegative dimensionless integer state and is not
defined for rates. Mixed modes or unequal simultaneous set values
reject. Unmentioned fields retain their pre-event values. State and
rate maps are handled separately, so a new rate does not retroactively
change the interval just realized.

World edits and all hinge edits are computed on clones. Only after
every check succeeds are the successor epoch, fired relation IDs,
hinge states, transition record, and accepted certificate published
together. No failed candidate partially mutates the world. The successor
at $\tau_*<H$ starts a new affine interval using $x^+,v^+$; the successor
at $H$ is terminal.

\subsection{Semantic replay, finite progress, and source corrections}

A hash certifies bytes, not that an asserted event follows from an epoch.
The adapter regenerates the expected certificate $C(E)$ from the current
immutable epoch and admits a submitted certificate $\widehat C$ only when
\begin{equation}
 \operatorname{canonical}(\widehat C)
       =\operatorname{canonical}(C(E)).
 \label{eq:r16-semantic-admission}
\end{equation}
This checks the root, source epoch, all candidates, complete simultaneous
membership, and all gate-derived claims. Altering a root or omitting an
earlier/tied relation and calculating a fresh hash does not satisfy this
equality. The source profile supplies the update bodies, so a certificate
cannot substitute different state edits.

An identical already accepted certificate returns a duplicate result
without a second transition. The same identity with changed bytes rejects.
Packed session loading reconstructs the initial profile, replays each
accepted certificate semantically, and compares the derived final epoch
and transitions; unpacking bytes alone is not accepted replay evidence.

Each successful event time is strictly later than the current start,
and at least one previously unconsumed relation is marked fired. For $m$
initial once-only relations there are at most $m$ event sets, followed
if necessary by a single event-free horizon finalization. A configured
event-set budget may stop earlier with an explicit resource result;
unresolved relations and conflicting edits stop without commit. This
argument does not extend to arbitrary rearming, newly created relations,
or zero-time cascades.

Three source implementation issues are corrected at this boundary.
The reviewed path accepts $\tau=H$ but then tries to create a forbidden
zero-length OpenSegment; R16 uses a terminal snapshot. The reviewed
rational dictionary parser can truncate floating limbs when its schema
has not been run; R16 requires actual integer limbs. The reviewed event
application can accept a rehashed claim without replaying its full
semantics; R16 recomputes the complete expected event set. These concrete
corrections strengthen the finite event contract. They do not add a
claim of arbitrary world prediction, a GPU analytic solver, or a new
physical conservation law.
